%% theory/eigen/counting.tex -- Task 2: explicit counting function and tail bounds.
%% Verification: theory/eigen/counting.py (pointwise bounds by quadrature; Area >= pi/21 and
%% n + 4g <= Area/pi + 4 on 4.3 million signatures; counting, tail and Hyp bounds on the committed
%% spectra of O(2,8,8) and O(3,3,12)).

\subsection{Counting eigenvalues with the trace formula}\label{sec:eig-count}

Write $0=\lambda_0<\lambda_1\le\lambda_2\le\dots$ for the eigenvalues of $\Orb$ with multiplicity, $\mathcal N_\Orb(x)=\#\{j:\lambda_j\le x\}$, and $\mathrm I,\mathrm E_m,\Hyp$ for the terms of Theorem~\ref{thm:IEH}.

\begin{lemma}[Elementary bounds]\label{eig:elem}
\leavevmode
\begin{enumerate}[label=(\roman*)]
\item Every $\Orb\in\Sig$ has $\Area(\Orb)\ge\pi/21$, and $n+4g\le\Area(\Orb)/\pi+4$.
\item For $t>0$ and $m\ge2$: $0<\mathrm E_m(t)\le b_0(m)=\frac{m^2-1}{12m}$ and $0<\mathrm I(t)\le e^{-t/4}\,\frac{\Area(\Orb)}{4\pi t}$.
\end{enumerate}
\end{lemma}

\begin{proof}
(i) Put $s=\Area/2\pi=2g-2+\sum_i(1-\frac1{m_i})>0$. If $g\ge2$, $s\ge2$; if $g=1$, $n\ge1$ and $s\ge\frac12$; if $g=0$ and $n\ge5$, $s\ge\frac12$. For $g=0$, $n=4$, the orders are not all $2$, so $s\ge2-\frac32-\frac13=\frac16$. For $g=0$, $n=3$, $p\le q\le r$: if $p\ge3$ then $r\ge4$ and $s\ge1-\frac13-\frac13-\frac14=\frac1{12}$; if $p=2$, then $q\ge3$, and $q=3$ forces $r\ge7$, $s\ge\frac1{42}$, $q=4$ forces $r\ge5$, $s\ge\frac1{20}$, and $q\ge5$ gives $s\ge\frac1{10}$. The second statement is the first line of the proof of Corollary~\ref{cor:sigarea}.
(ii) $0<h_t(r)\le1$, and the integrands of Theorem~\ref{thm:IEH} are positive. By Euler's beta integral (Section~\sref{sec:derivation} of the supplement, at $s=0$), $\int_\R e^{-2\theta_jr}(1+e^{-2\pi r})^{-1}dr=(2\sin\theta_j)^{-1}$, so $\mathrm E_m(t)\le\sum_j(4m\sin^2\theta_j)^{-1}=\Phi_m(0)=\phi_0(m)=b_0(m)$ (Lemma~\ref{lem:Phi}, \eqref{eq:phik}). And $0\le r\tanh\pi r\le|r|$, $\int_\R|r|e^{-tr^2}dr=1/t$.
\end{proof}

\begin{proposition}[Counting and tails through the heat trace]\label{eig:counttail}
For every $\Orb$, $x>0$, $\Lambda>0$ and $0<s<t$:
\[
  \mathcal N_\Orb(x)\le e\,\cZ_\Orb(1/x),\qquad \sum_{\lambda_j>\Lambda}e^{-\lambda_jt}\le e^{-\Lambda(t-s)}\,\cZ_\Orb(s).
\]
If moreover $N$ is an integer with $N\ge e\,\cZ_\Orb(1/\Lambda)$, then $\lambda_j>\Lambda$ for every $j\ge N$, so $\sum_{j\ge N}e^{-\lambda_jt}\le e^{-\Lambda(t-s)}\cZ_\Orb(s)$.
\end{proposition}

\begin{proof}
$\cZ(1/x)\ge\sum_{\lambda_j\le x}e^{-\lambda_j/x}\ge e^{-1}\mathcal N(x)$, and $e^{-\lambda t}=e^{-\lambda(t-s)}e^{-\lambda s}\le e^{-\Lambda(t-s)}e^{-\lambda s}$ for $\lambda>\Lambda$. If $N\ge e\cZ(1/\Lambda)\ge\mathcal N(\Lambda)$, the indices $j$ with $\lambda_j\le\Lambda$ are $0,\dots,\mathcal N(\Lambda)-1<N$.
\end{proof}

Lemma~\ref{lem:hypbound} bounds $\Hyp(t)$ only for $t\le\ell^2/(2(1+\ell))$. For all $t$:

\begin{proposition}[The hyperbolic term at every time]\label{eig:hypall}
For every $\Orb\in\Sig$ with systole $\ell$ and every $t>0$,
\[
  0<\Hyp(t)\ \le\ \frac{2\sqrt\pi\,e^{3\diam}}{\Area(\Orb)\,(1-e^{-\ell})}\,e^{\frac{17t}4-\frac12-\ell}.
\]
\end{proposition}

\begin{proof}
As in the proof of Lemma~\ref{lem:hypbound}, $\ell(\gamma_0)\le\ell(\gamma)$ and $2\sinh\frac x2\ge e^{x/2}(1-e^{-\ell})$ for $x\ge\ell$ give $\Hyp(t)\le e^{-t/4}((1-e^{-\ell})\sqrt{4\pi t})^{-1}\sum_{[\gamma]}\ell(\gamma)e^{-\ell(\gamma)/2-\ell(\gamma)^2/4t}$. For $x\ge0$,
\[
  xe^{-\frac x2-\frac{x^2}{4t}}=\big(xe^{-\frac{x^2}{8t}}\big)\big(e^{\frac{3x}2-\frac{x^2}{8t}}\big)e^{-2x}\le2\sqrt t\,e^{-\frac12}\,e^{\frac{9t}2}\,e^{-2x},
\]
the two maxima being attained at $x=2\sqrt t$ and $x=6t$. By Lemma~\ref{lem:counting}, $n_\Orb(x)\le\pi e^{x+3\diam}/\Area$, and integration by parts ($n_\Orb=0$ on $[0,\ell)$, $n_\Orb(X)e^{-2X}\to0$) gives $\sum_{[\gamma]}e^{-2\ell(\gamma)}=2\int_\ell^\infty n_\Orb(x)e^{-2x}dx\le2\pi e^{3\diam-\ell}/\Area$. Multiply.
\end{proof}

\begin{theorem}[Explicit counting and tail bounds]\label{eig:count}
Let $\Orb\in\Cl(A,\varepsilon,M)$, $D=D(A,\varepsilon,M)$ (Theorem~\ref{eig:diam}), $n_*=\lfloor A/\pi\rfloor+4$, $\beta_M=n_*\frac{M^2-1}{12M}$, and
\[
  \mathcal H(t)=\begin{cases}\dfrac{21\,e^{3D}\,\varepsilon e^{\varepsilon/2}}{1-e^{-\varepsilon}}\Big(1+\dfrac{2t}{\varepsilon-t}\Big)\dfrac{e^{-\varepsilon^2/4t}}{\sqrt{4\pi t}}, & 0<t\le\dfrac{\varepsilon^2}{2(1+\varepsilon)},\\[3mm]
  \dfrac{42\,e^{3D}}{\sqrt\pi\,(1-e^{-\varepsilon})}\,e^{\frac{17t}4-\frac12-\varepsilon}, & t>\dfrac{\varepsilon^2}{2(1+\varepsilon)}.\end{cases}
\]
Then for all $t>0$,
\[
  \cZ_\Orb(t)\ \le\ \cZ_b(t):=\frac A{4\pi t}+\beta_M+\mathcal H(t),
\]
so $\mathcal N_\Orb(x)\le e\,\cZ_b(1/x)$ for $x>0$, and $\sum_{j\ge N}e^{-\lambda_jt}\le e^{-\Lambda(t-s)}\cZ_b(s)$ whenever $0<s<t$ and $N\ge e\,\cZ_b(1/\Lambda)$.
\end{theorem}

\begin{proof}
Theorem~\ref{thm:IEH}, Lemma~\ref{eig:elem} ($n\le n_*$, and $b_0$ is increasing), and for $\Hyp$: Lemma~\ref{lem:hypbound} or Proposition~\ref{eig:hypall} with $\diam<D$, $\Area\ge\pi/21$, and $\ell\ge\varepsilon$. On $t\le\varepsilon^2/(2(1+\varepsilon))$ the bound $B$ of Lemma~\ref{lem:hypbound} decreases in $\ell$ along $[\varepsilon,\ell]$ (as in the proof of Theorem~\ref{thm:quantlocality}(a)), and the bound of Proposition~\ref{eig:hypall} decreases in $\ell$. Then Proposition~\ref{eig:counttail}.
\end{proof}

So $\mathcal N_\Orb(x)\le\frac{e\,A}{4\pi}x+e\,\beta_M+e\,\mathcal H(1/x)$, with $\mathcal H(1/x)\to0$ as $x\to\infty$. On the committed spectra of $\Orb(2,8,8)$ and $\Orb(3,3,12)$ (about $2850$ eigenvalues each), $\mathcal N(x)/(e\,\cZ_b(1/x))\le0.37$ for $5\le x\le15000$.
